\begin{figure}[H]
    \centering
\begin{tikzpicture}[
    scale=1.5,
    cota/.style={<->, >=stealth, thin, gray!80!black},  % Estilo das linhas de cota
    retangulo/.style={draw=gray!80!black, thick},        % Estilo do retângulo
    divisao/.style={draw=black, thick, dashed},          % Linhas de divisão (charneiras)
    eixo/.style={->, >=stealth, thick, gray!80!black},
    corte/.style={draw=black, dash dot, thick},          % Linhas dos cortes A e B
    setaCorte/.style={->, >=stealth, thick, black}
]

    % 1. Coordenadas do retângulo
    \coordinate (A) at (0, 0);
    \coordinate (B) at (2, 0);
    \coordinate (C) at (2, 4);
    \coordinate (D) at (0, 4);

    % Interseção P da charneira (ponto de encontro dos ângulos de 60°)
    % Forma um triângulo à direita e um trapézio à esquerda
    \coordinate (P1) at (0, 1.5);
    \coordinate (P2) at (0, 2.5);

    % 2. Hachuras externas no topo e na direita (sem borda)
    %\path[pattern={Lines[angle=45, distance=6pt, line width=0.5pt]}, pattern color=gray!60] (0, 3) rectangle (2.3, 3.3);
    \path[pattern={Lines[angle=45, distance=6pt, line width=0.5pt]}, pattern color=gray!60] (2, 0) rectangle (2.3, 4);

    % 3. Desenho do Retângulo
    \draw[retangulo] (A) rectangle (C);

    % 4. Divisão em Triângulo (direita) e Trapézio (esquerda)
    \draw[divisao] (C) -- (P2) --(P1) -- (B);

    % 5. Linhas dos Cortes A (Horizontal) e B (Vertical)
    % Corte A (Horizontal em y = 1.5)
    \draw[corte] (-0.4, 2) -- (2.5, 2);
    \draw[setaCorte] (-0.4, 2) -- (-0.4, 2.25) node[above, font=\small\bfseries] {A};
    \draw[setaCorte] (2.5, 2) -- (2.5, 2.25) node[above, font=\small\bfseries] {A};

    % Corte B (Vertical em x = 1.0)
    \draw[corte] (1.0, -0.6) -- (1.0, 4.5);
    \draw[setaCorte] (1.0, 4.5) -- (1.3,4.5) node[right, font=\small\bfseries] {B};
    \draw[setaCorte] (1.0, -0.6) -- (1.3, -0.6) node[right, font=\small\bfseries] {B};
    
    \node[font=\small\bfseries, gray!90!black, fill=white, inner sep=1.5pt, circle] at (0.5, 3.5) {$A_y$};
    %\node[font=\small\bfseries, gray!90!black, fill=white, inner sep=1.5pt, circle] at (1.23, 1.6) {$A_x$};
    \node[font=\small\bfseries, gray!90!black, fill=white, inner sep=1.5pt, circle] at (0.5, 0.5) {$A_y$};
    \node[font=\small\bfseries, gray!90!black, fill=white, inner sep=1.5pt, circle] at (1, 2) {$A'_x$};

    % 6. Indicações dos Ângulos de 60° na borda direita
    % Canto superior direito (C): 60° em relação à borda vertical
    \draw[thin, gray!80!black] (2, 3.3) arc (270:210:0.7);
    \node[font=\tiny, gray!90!black] at (1.75, 3.2) {$60^\circ$};

    % Canto inferior direito (B): 60° em relação à borda vertical
    \draw[thin, gray!80!black] (2, 0.7) arc (90:150:0.7);
    \node[font=\tiny, gray!90!black] at (1.75, 0.8) {$60^\circ$};

    % 7. Eixos de Coordenadas x e y
    \draw[eixo] (0.2, 0.2) -- (0.7, 0.2) node[above=-2pt, font=\tiny\bfseries] {$x$};
    \draw[eixo] (0.2, 0.2) -- (0.2, 0.7) node[right=-2pt, font=\tiny\bfseries] {$y$};

    % 8. Cotas de Dimensão
    %\def\offsetH{0.4}
    %\draw[cota] (0, -\offsetH) -- (2, -\offsetH) 
        %node[midway, below=3pt, font=\small] {$476\text{ cm}$};

    %\def\offsetV{0.4}
    %\draw[cota] (-\offsetV, 0) -- (-\offsetV, 3) 
        %node[midway, left=3pt, font=\small, rotate=90, anchor=south] {$556 \text{ cm}$};

\end{tikzpicture}
    %\caption{Retângulo dividido em triângulo e trapézio com cortes A e B}
    %\label{fig:retangulo_triangulo_trapezio}
\end{figure}